\documentclass{beamer}
\usetheme{Madrid}

\usepackage[utf8]{inputenc} % solo si usas pdflatex
\usepackage[T1]{fontenc}    % solo si usas pdflatex

\usepackage{amsmath,amssymb}
\usepackage{microtype}
\usepackage{cancel}         % si lo usas
\usepackage{fontawesome}    % si lo usas
% \usepackage{xypic}        % solo si lo usas
% \usepackage{multicol}     % normalmente no; usar columns de beamer
% \usepackage{amsthm}       % solo si defines teoremas propios
% \usepackage{scrextend}    % solo si lo usas

\usepackage{graphicx}
\graphicspath{{images/}}

\hypersetup{pdfborder={0 0 0}}


\title{Introducción a Azure}
\author{Rodrigo Blanco Betes. Alfa Formación}
\date{Enero del 2026}

\begin{document}

\begin{frame}
\titlepage
\end{frame}

\begin{frame}
\tableofcontents
\end{frame}

\section{Introducción a Azure}

\begin{frame}{Infraestructura en Azure}
\begin{itemize}
    \item Organización de recursos y servicios en la nube
    \item Gestión de costes, entornos y despliegues
    \item Recursos principales:
    \begin{itemize}
        \item Cuentas
        \item Resource Groups
        \item Virtual Machines
        \item AKS
    \end{itemize}
\end{itemize}
\end{frame}

\begin{frame}{Cuentas en Azure}
\begin{itemize}
    \item Punto de partida para consumir servicios en Azure
    \item Asociadas a suscripción, facturación y control de acceso
    \item Permiten:
    \begin{itemize}
        \item Monitorizar costes
        \item Asignar presupuestos
        \item Controlar el uso por equipos o proyectos
    \end{itemize}
    \item Base para la gobernanza y la trazabilidad del gasto
\end{itemize}
\end{frame}

\begin{frame}{Resource Groups}
\begin{itemize}
    \item Contenedor lógico para agrupar recursos relacionados
    \item Facilitan la organización según distintos criterios:
    \begin{itemize}
        \item Entornos: desarrollo, pruebas, producción
        \item Funcionalidad: frontend, backend, datos
        \item Área o centro de costes
    \end{itemize}
    \item Simplifican la administración, permisos y eliminación conjunta
\end{itemize}
\end{frame}

\begin{frame}{Virtual Machines}
\begin{itemize}
    \item Servicio de computación IaaS en Azure
    \item Permite desplegar servidores con control completo del sistema operativo
    \item Casos de uso habituales:
    \begin{itemize}
        \item Aplicaciones legacy
        \item Entornos de laboratorio
        \item Servicios que requieren configuración específica
    \end{itemize}
\end{itemize}
\end{frame}

\begin{frame}{AKS}
\begin{itemize}
    \item Azure Kubernetes Service
    \item Plataforma gestionada para ejecutar contenedores orquestados con Kubernetes
    \item Reduce la complejidad operativa del clúster
    \item Adecuado para:
    \begin{itemize}
        \item Aplicaciones distribuidas
        \item Microservicios
        \item Despliegues escalables y resilientes
    \end{itemize}
\end{itemize}
\end{frame}

\section{Primeras prácticas}

\begin{frame}{Gestión de cuentas y costes}
\begin{itemize}
    \item Identificación de suscripciones y responsables
    \item Revisión del consumo y control presupuestario
    \item Buenas prácticas:
    \begin{itemize}
        \item Etiquetado de recursos
        \item Separación por proyectos o departamentos
        \item Seguimiento periódico del gasto
    \end{itemize}
\end{itemize}
\end{frame}

\begin{frame}{Uso práctico de Resource Groups}
\begin{itemize}
    \item Crear grupos de recursos alineados con la estructura del proyecto
    \item Separar recursos por entorno:
    \begin{itemize}
        \item \texttt{rg-dev}
        \item \texttt{rg-test}
        \item \texttt{rg-prod}
    \end{itemize}
    \item Aplicar nomenclatura consistente
    \item Facilitar mantenimiento y limpieza de recursos
\end{itemize}
\end{frame}

\begin{frame}{Virtual Machine: creación y configuración}
\begin{itemize}
    \item Selección de imagen, tamaño y red
    \item Configuración inicial:
    \begin{itemize}
        \item Sistema operativo
        \item Credenciales de acceso
        \item Reglas de red y puertos
        \item Discos y monitorización
    \end{itemize}
    \item Validación de conectividad y estado operativo
\end{itemize}
\end{frame}

\begin{frame}{Virtual Machine: decomisionado}
\begin{itemize}
    \item Apagado y revisión previa del recurso
    \item Copia o respaldo de datos necesarios
    \item Eliminación controlada de:
    \begin{itemize}
        \item Máquina virtual
        \item Discos asociados
        \item Interfaces de red e IP públicas
    \end{itemize}
    \item Objetivo: evitar costes residuales y recursos huérfanos
\end{itemize}
\end{frame}

\begin{frame}{AKS: creación y configuración}
\begin{itemize}
    \item Creación del clúster y definición de nodos
    \item Configuración de red, identidad y escalado
    \item Despliegue inicial de cargas de trabajo en contenedores
    \item Revisión de:
    \begin{itemize}
        \item Acceso al clúster
        \item Estado de nodos y pods
        \item Integración con monitorización
    \end{itemize}
\end{itemize}
\end{frame}

\begin{frame}{AKS: decomisionado}
\begin{itemize}
    \item Revisión de servicios y cargas desplegadas
    \item Exportación o migración de configuraciones necesarias
    \item Eliminación del clúster y recursos asociados
    \item Verificación final para evitar consumos innecesarios
\end{itemize}
\end{frame}

\begin{frame}{Resumen}
\begin{itemize}
    \item Azure permite organizar infraestructura de forma flexible y gobernada
    \item Los Resource Groups ayudan a estructurar recursos y costes
    \item Las Virtual Machines ofrecen control total sobre entornos IaaS
    \item AKS facilita la ejecución y operación de contenedores a escala
    \item La creación, configuración y decomisionado forman parte del ciclo de vida operativo
\end{itemize}
\end{frame}

\section{Procesamiento de datos}

\begin{frame}{Arquitectura de procesamiento}
\begin{itemize}
    \item Flujo de datos estructurado en varias fases
    \item Separación por capas para facilitar mantenimiento y trazabilidad
    \item Etapas principales del procesamiento:
    \begin{itemize}
        \item Extracción
        \item Ingesta
        \item Limpieza
        \item Generación de tablones
        \item Construcción del modelo de explotación
    \end{itemize}
\end{itemize}
\end{frame}

\begin{frame}{Extracciones: origen a Blob}
\begin{itemize}
    \item Obtención de datos desde los sistemas de origen
    \item Posibles fuentes:
    \begin{itemize}
        \item Bases de datos
        \item APIs
        \item Sistemas externos
        \item Ficheros
    \end{itemize}
    \item Los datos se almacenan inicialmente en almacenamiento Blob
    \item Objetivo: disponer de una copia inicial sin transformaciones
\end{itemize}
\end{frame}

\begin{frame}{Ingesta: Blob a Staging}
\begin{itemize}
    \item Movimiento de datos desde Blob hacia la capa \textit{staging}
    \item Zona intermedia para preparar el procesamiento
    \item Características principales:
    \begin{itemize}
        \item Datos cercanos al formato original
        \item Preparados para su transformación
        \item Control de versiones o particiones
    \end{itemize}
\end{itemize}
\end{frame}

\begin{frame}{Limpieza con Spark: Staging a Golden}
\begin{itemize}
    \item Procesamiento de datos mediante Apache Spark
    \item Aplicación de transformaciones de calidad:
    \begin{itemize}
        \item Eliminación de registros inválidos
        \item Normalización de formatos
        \item Tratamiento de valores nulos
        \item Validaciones de integridad
    \end{itemize}
    \item Resultado almacenado en la capa \textit{golden}
\end{itemize}
\end{frame}

\begin{frame}{Generación de tablones con Spark}
\begin{itemize}
    \item Construcción de datasets agregados o enriquecidos
    \item Combinación de distintas fuentes de datos
    \item Operaciones habituales:
    \begin{itemize}
        \item Joins
        \item Agregaciones
        \item Cálculo de métricas
    \end{itemize}
    \item Los resultados se almacenan en la capa \textit{golden}
\end{itemize}
\end{frame}

\begin{frame}{Modelo de explotación: Golden a Silver}
\begin{itemize}
    \item Generación de estructuras optimizadas para análisis
    \item Procesamiento mediante Spark
    \item Organización orientada a consumo analítico
    \item Características:
    \begin{itemize}
        \item Datos consolidados
        \item Estructuras optimizadas para consultas
        \item Preparación para herramientas de BI o analítica
    \end{itemize}
\end{itemize}
\end{frame}

\begin{frame}{Modelo de explotación: Golden a Cosmos}
\begin{itemize}
    \item Publicación de datos hacia Azure Cosmos DB
    \item Orientado a aplicaciones operacionales o APIs
    \item Ventajas:
    \begin{itemize}
        \item Baja latencia
        \item Alta escalabilidad
        \item Acceso directo desde aplicaciones
    \end{itemize}
    \item Los datos proceden de la capa \textit{golden}
\end{itemize}
\end{frame}

\begin{frame}{Resumen del flujo de procesamiento}
\begin{itemize}
    \item Extracción de datos desde sistemas origen
    \item Ingesta inicial en almacenamiento intermedio
    \item Procesamiento y limpieza mediante Spark
    \item Generación de datasets enriquecidos
    \item Publicación de modelos para analítica y aplicaciones
\end{itemize}
\end{frame}

\section{Visualización}

\begin{frame}{Visualización en Azure}
\begin{itemize}
    \item La visualización permite analizar el estado de los sistemas y explotar los datos procesados
    \item En entornos Azure, las herramientas de visualización suelen cubrir:
    \begin{itemize}
        \item Monitorización operativa
        \item Observabilidad de infraestructura y aplicaciones
        \item Exploración y análisis de datos
    \end{itemize}
    \item Herramientas principales:
    \begin{itemize}
        \item Kibana
        \item Grafana
        \item Notebooks
    \end{itemize}
\end{itemize}
\end{frame}

\begin{frame}{Kibana en entornos Azure}
\begin{itemize}
    \item Herramienta orientada a exploración y visualización de datos indexados
    \item Uso habitual junto con Elasticsearch
    \item Permite:
    \begin{itemize}
        \item Crear dashboards
        \item Analizar logs
        \item Investigar incidencias
        \item Realizar búsquedas sobre grandes volúmenes de información
    \end{itemize}
    \item En Azure puede integrarse en arquitecturas de observabilidad y análisis de logs
\end{itemize}
\end{frame}

\begin{frame}{Grafana en Azure}
\begin{itemize}
    \item Plataforma de visualización orientada a métricas y monitorización
    \item Muy útil para cuadros de mando técnicos y operativos
    \item Permite:
    \begin{itemize}
        \item Visualizar series temporales
        \item Supervisar infraestructura
        \item Monitorizar aplicaciones y pipelines
        \item Centralizar indicadores de distintos orígenes
    \end{itemize}
    \item En Azure puede conectarse con servicios de monitorización y fuentes de datos diversas
\end{itemize}
\end{frame}

\begin{frame}{Notebooks en Azure}
\begin{itemize}
    \item Entorno interactivo para exploración, transformación y análisis de datos
    \item Permiten combinar:
    \begin{itemize}
        \item Código
        \item Resultados
        \item Visualizaciones
        \item Documentación
    \end{itemize}
    \item Uso habitual en Azure para:
    \begin{itemize}
        \item Análisis exploratorio
        \item Desarrollo con Spark
        \item Validación de datasets
        \item Prototipado de modelos
    \end{itemize}
\end{itemize}
\end{frame}

\begin{frame}{Resumen de visualización}
\begin{itemize}
    \item Kibana facilita la exploración de logs y eventos
    \item Grafana está orientado a métricas y monitorización
    \item Los notebooks permiten análisis interactivo y desarrollo de procesos de datos
    \item En Azure, estas herramientas complementan la explotación y la observabilidad de la plataforma
\end{itemize}
\end{frame}

\section{Almacenamiento}

\begin{frame}{Almacenamiento en Azure}
\begin{itemize}
    \item El almacenamiento es una pieza central en cualquier arquitectura de datos
    \item Azure ofrece distintos modelos según el tipo de información y su uso
    \item Opciones principales:
    \begin{itemize}
        \item Modelo SQL
        \item Data Lake
        \item Storage
        \item Cosmos DB
    \end{itemize}
    \item Cada una responde a necesidades diferentes de estructura, escalabilidad y acceso
\end{itemize}
\end{frame}

\begin{frame}{Modelo SQL}
\begin{itemize}
    \item Almacenamiento relacional orientado a datos estructurados
    \item Basado en tablas, relaciones, claves e integridad referencial
    \item Adecuado para:
    \begin{itemize}
        \item Aplicaciones transaccionales
        \item Modelos normalizados
        \item Consultas SQL
        \item Reporting estructurado
    \end{itemize}
    \item En Azure suele emplearse cuando se necesita consistencia y esquema definido
\end{itemize}
\end{frame}

\begin{frame}{Data Lake en Azure}
\begin{itemize}
    \item Repositorio escalable para grandes volúmenes de datos
    \item Admite datos estructurados, semiestructurados y no estructurados
    \item Organización habitual por capas:
    \begin{itemize}
        \item \textbf{Landing}: recepción inicial de datos
        \item \textbf{Golden}: datos limpios, validados y enriquecidos
        \item \textbf{Silver}: datos preparados para explotación o consumo específico
    \end{itemize}
    \item Base habitual para arquitecturas analíticas modernas en Azure
\end{itemize}
\end{frame}

\begin{frame}{Capa Landing}
\begin{itemize}
    \item Zona de entrada de datos procedentes de sistemas origen
    \item Conserva la información con mínima transformación
    \item Objetivos:
    \begin{itemize}
        \item Mantener trazabilidad
        \item Preservar el dato original
        \item Facilitar reprocesados
    \end{itemize}
    \item Es la primera parada en el flujo de ingesta dentro del Data Lake
\end{itemize}
\end{frame}

\begin{frame}{Capas Golden y Silver}
\begin{itemize}
    \item \textbf{Golden}:
    \begin{itemize}
        \item Datos depurados y consolidados
        \item Resultado de procesos de limpieza y enriquecimiento
    \end{itemize}
    \item \textbf{Silver}:
    \begin{itemize}
        \item Datos orientados a consumo
        \item Preparados para análisis, explotación o integración con otros sistemas
    \end{itemize}
    \item Estas capas ayudan a separar el dato bruto del dato de negocio
\end{itemize}
\end{frame}

\begin{frame}{Storage en Azure: Blob y File}
\begin{itemize}
    \item Azure Storage proporciona distintos servicios según el tipo de acceso requerido
    \item \textbf{Blob Storage}:
    \begin{itemize}
        \item Orientado a objetos
        \item Adecuado para ficheros, datos masivos, backups y datasets
    \end{itemize}
    \item \textbf{File Storage}:
    \begin{itemize}
        \item Comparte archivos mediante estructura tipo sistema de ficheros
        \item Útil para recursos compartidos entre aplicaciones o servidores
    \end{itemize}
\end{itemize}
\end{frame}

\begin{frame}{Blob Storage}
\begin{itemize}
    \item Servicio de almacenamiento de objetos en Azure
    \item Muy utilizado en arquitecturas de datos y analítica
    \item Casos de uso:
    \begin{itemize}
        \item Ingesta de datos
        \item Almacenamiento de ficheros de entrada y salida
        \item Persistencia de datos intermedios
        \item Integración con pipelines y procesos Spark
    \end{itemize}
\end{itemize}
\end{frame}

\begin{frame}{File Storage}
\begin{itemize}
    \item Servicio orientado a compartición de archivos
    \item Permite montar recursos compartidos en entornos Windows o Linux
    \item Adecuado para:
    \begin{itemize}
        \item Carpetas compartidas
        \item Aplicaciones legacy
        \item Intercambio de ficheros entre servicios
    \end{itemize}
    \item Complementa al almacenamiento tipo Blob cuando se necesita acceso tipo fichero
\end{itemize}
\end{frame}

\begin{frame}{Cosmos DB en Azure}
\begin{itemize}
    \item Base de datos NoSQL distribuida globalmente
    \item Diseñada para alta disponibilidad, baja latencia y escalabilidad horizontal
    \item Modelos mencionados:
    \begin{itemize}
        \item API SQL
        \item API MongoDB
    \end{itemize}
    \item Adecuada para aplicaciones con necesidades de acceso rápido y flexible a los datos
\end{itemize}
\end{frame}

\begin{frame}{Cosmos DB: SQL y MongoDB}
\begin{itemize}
    \item \textbf{API SQL}:
    \begin{itemize}
        \item Modelo documental nativo
        \item Consultas sobre documentos JSON
        \item Muy utilizada en desarrollos cloud nativos
    \end{itemize}
    \item \textbf{API MongoDB}:
    \begin{itemize}
        \item Compatible con el ecosistema y herramientas de MongoDB
        \item Facilita migraciones o integraciones con aplicaciones existentes
    \end{itemize}
    \item Ambas permiten trabajar con datos semiestructurados en Azure
\end{itemize}
\end{frame}

\begin{frame}{Resumen de almacenamiento}
\begin{itemize}
    \item El modelo SQL cubre necesidades relacionales y estructuradas
    \item El Data Lake organiza el ciclo de vida del dato en distintas capas
    \item Blob y File resuelven necesidades diferentes de almacenamiento en Azure
    \item Cosmos DB ofrece una alternativa NoSQL escalable para aplicaciones modernas
\end{itemize}
\end{frame}

\section{Procesamiento en tiempo real}

\begin{frame}{Arquitectura de datos en tiempo real}
\begin{itemize}
    \item El procesamiento en tiempo real permite analizar datos según se generan
    \item Se utiliza en escenarios donde la latencia debe ser mínima
    \item Ejemplos de uso:
    \begin{itemize}
        \item IoT y sensores
        \item Monitorización de sistemas
        \item Procesamiento de eventos
        \item Analítica operativa
    \end{itemize}
    \item En Azure, el flujo suele incluir ingestión de eventos, procesamiento streaming y almacenamiento
\end{itemize}
\end{frame}

\begin{frame}{Node-RED}
\begin{itemize}
    \item Herramienta de programación visual basada en flujos
    \item Permite integrar dispositivos, APIs y servicios cloud
    \item Uso habitual en entornos IoT:
    \begin{itemize}
        \item Conectar sensores o dispositivos
        \item Transformar eventos
        \item Enviar datos a servicios de Azure
    \end{itemize}
    \item Facilita la creación rápida de pipelines de eventos
\end{itemize}
\end{frame}

\begin{frame}{Azure IoT Hub}
\begin{itemize}
    \item Servicio gestionado para comunicación con dispositivos IoT
    \item Permite conectar y gestionar miles o millones de dispositivos
    \item Funcionalidades principales:
    \begin{itemize}
        \item Ingesta de telemetría
        \item Gestión de dispositivos
        \item Comunicación bidireccional
        \item Seguridad y autenticación
    \end{itemize}
    \item Punto de entrada habitual para datos procedentes de sensores
\end{itemize}
\end{frame}

\begin{frame}{Azure Event Hub}
\begin{itemize}
    \item Plataforma de ingestión de eventos a gran escala
    \item Diseñada para manejar grandes volúmenes de datos en streaming
    \item Características principales:
    \begin{itemize}
        \item Alta capacidad de ingestión
        \item Procesamiento distribuido
        \item Integración con servicios de análisis
    \end{itemize}
    \item Muy utilizada como buffer entre productores de eventos y sistemas de procesamiento
\end{itemize}
\end{frame}

\begin{frame}{Spark Streaming}
\begin{itemize}
    \item Extensión de Apache Spark para procesamiento de datos en streaming
    \item Permite analizar flujos de eventos en tiempo casi real
    \item Funcionalidades:
    \begin{itemize}
        \item Transformación de datos
        \item Agregaciones en streaming
        \item Cálculo de métricas en tiempo real
    \end{itemize}
    \item Puede consumir eventos desde Event Hub u otras fuentes
\end{itemize}
\end{frame}

\begin{frame}{Almacenamiento en Cosmos DB}
\begin{itemize}
    \item Base de datos NoSQL distribuida utilizada para persistir resultados
    \item Adecuada para escenarios de baja latencia
    \item En arquitecturas streaming permite:
    \begin{itemize}
        \item Guardar resultados procesados
        \item Servir datos a aplicaciones en tiempo real
        \item Escalar horizontalmente con gran volumen de eventos
    \end{itemize}
\end{itemize}
\end{frame}

\begin{frame}{Notebooks en procesamiento streaming}
\begin{itemize}
    \item Permiten desarrollar y validar pipelines de procesamiento
    \item Usados habitualmente en entornos Spark
    \item Utilidades:
    \begin{itemize}
        \item Desarrollo de lógica de transformación
        \item Análisis exploratorio de flujos de datos
        \item Monitorización y validación de resultados
    \end{itemize}
\end{itemize}
\end{frame}

\begin{frame}{Flujo típico de procesamiento en tiempo real}
\begin{itemize}
    \item Generación de eventos desde dispositivos o aplicaciones
    \item Ingesta mediante IoT Hub o Event Hub
    \item Procesamiento en streaming con Spark
    \item Persistencia de resultados en sistemas de almacenamiento
    \item Visualización o explotación de la información
\end{itemize}
\end{frame}

\section{Gobierno del dato}

\begin{frame}{Gobierno del dato}
\begin{itemize}
    \item El gobierno del dato define cómo se gestionan los datos dentro de una organización
    \item Su objetivo es garantizar:
    \begin{itemize}
        \item Calidad del dato
        \item Seguridad
        \item Cumplimiento normativo
        \item Uso correcto de la información
    \end{itemize}
\end{itemize}
\end{frame}

\begin{frame}{Elementos clave del gobierno del dato}
\begin{itemize}
    \item Definición de responsables y roles sobre los datos
    \item Gestión de la calidad y validación de información
    \item Control de accesos y seguridad
    \item Catalogación y documentación de datasets
    \item Seguimiento del ciclo de vida del dato
\end{itemize}
\end{frame}

\begin{frame}{Gobierno del dato en entornos Azure}
\begin{itemize}
    \item Azure proporciona herramientas para gestionar el ciclo de vida del dato
    \item Permite:
    \begin{itemize}
        \item Controlar accesos a los recursos
        \item Auditar el uso de los datos
        \item Catalogar y descubrir información
        \item Aplicar políticas de seguridad
    \end{itemize}
    \item El gobierno del dato es esencial en plataformas de datos empresariales
\end{itemize}
\end{frame}

\begin{frame}{Resumen}
\begin{itemize}
    \item El procesamiento en tiempo real permite analizar eventos según se generan
    \item Azure ofrece servicios para ingestión, procesamiento y almacenamiento streaming
    \item Cosmos DB facilita el acceso rápido a los resultados
    \item El gobierno del dato garantiza calidad, seguridad y control sobre la información
\end{itemize}
\end{frame}

\end{document}
